\reportsection{Background}
\label{sec:background}

Traditionally, when we think of Merkle trees we think of arity $2$ throughout, one hasher
$f$ at every level, and a simple binding claim tied to the root commitment.  As protocols
mature and the area is studied more deeply, the picture has broadened to varying arity,
mixing hashers across levels, and different flavours of binding, each motivated by
performance and recursion-friendliness. That shift complicates security analysis, and it
is easy to lose track of details that a fixed arity-$2$, single-hasher tree used to
settle by construction.

\begin{definition}[Vector commitment scheme]\label{def:vcgeneral}
A \emph{vector commitment scheme}~\cite{CF13} over alphabet $\MTAlphabet$ and length
$\MTMessageLength$ is a triple $\VC=(\Commit,\Open,\VCCheck)$:
\begin{itemize}
  \item $\Commit(\msgvec)\to(\comm,\aux)$ binds a message
  $\msgvec\in\MTAlphabet^{\MTMessageLength}$ to a short commitment $\comm$ and
  retains a state $\aux$;
  \item $\Open(\aux,\idxset)\to\opening$ authenticates an index set
  $\idxset\subseteq[\MTMessageLength]$;
  \item $\VCCheck(\comm,\idxset,\MTMessageSubVector,\opening)\to\bits$ decides
  whether $\opening$ authenticates the subvector $\MTMessageSubVector$ at
  $\idxset$ against $\comm$, where for an honest committer
  $\MTMessageSubVector$ is exactly the committed message $\msgvec$ restricted
  to the positions in $\idxset$.
\end{itemize}
A scheme is \emph{complete} (\cref{def:vc-complete}) and \emph{position-binding}
(\cref{def:vc-binding}).  It may also implement optional extensions, namely hiding
(\cref{def:vc-hiding}), updatability, aggregation, functional opening, and
knowledge soundness, of which this report treats only hiding.
\end{definition}

\begin{definition}[Completeness]\label{def:vc-complete}
$\VC$ is \emph{complete} if for every $\msgvec\in\MTAlphabet^{\MTMessageLength}$
and every $\idxset\subseteq[\MTMessageLength]$ the honest
$\Commit$--$\Open$--$\VCCheck$ round trip accepts $\msgvec_{\idxset}$ with
probability~$1$.  For Merkle this is honest reconstruction of the root from an
opening.
\end{definition}

\begin{definition}[Merkle commitment scheme]
The \emph{Merkle commitment scheme}~\cite{Merkle88}
$\MTSymbol=(\MTCommit,\allowbreak\MTOpen,\allowbreak\MTCheck)$
instantiates \cref{def:vcgeneral} in the random-oracle model over a general
Merkle shape $G_\ell$ of depth $d=\max_i d_i$, the deepest of the per-leaf depths
$d_i$ over leaves $i$: the three algorithms query a
shared oracle $\ROFunction$, the commitment is the tree root, and the committer
state is the retained salts and labels.  The binary tree is the arity-$2$ subset
of this shape.  Two parameters carry the security: the digest
width $\digestbits$ governs binding and the salt size $\MTSaltSize$ governs
hiding, with no tuning of one supplying what the other lacks.
\end{definition}

Binding is required of every scheme in this report, and hiding matters only where the
protocol built on top needs zero-knowledge. At the \code{MerkleHasher} trait of
\textsf{ark-mt} the digest and salt widths are associated
types, fixed by whoever writes the implementation, and nothing compares them against the
size of the field.  BabyBear at $\secpar=128$ therefore gets a one-element digest over a
${\sim}30$-bit space and a one-element salt carrying ${\sim}30$ bits of entropy.  The
book~\cite{CY26} reaches the same figure and observes that such an instantiation is not
flawed as such, since a protocol may bound the prover's Merkle queries far below the
birthday threshold and argue security from its own soundness analysis.  The trait gives
no way to declare that bound, which is the role of $h$ in the rule below.

%% V2 copy of the design rule, boxed in Background so Binding and Hiding can
%% instantiate it without forward reference.  This box states the rule, and the
%% justification follows it in Background.

\begin{rulebox}
\noindent\textbf{The design rule.}\quad
A protocol designer supplies five things: a field~$\field$, a security
level~$\secpar$, a hash~$H$, a choice of standard or hiding interface, and an
adversary model, declared as the query-budget exponent~$h$.  Everything else
is derived.  From $(\field,\secpar,h)$ the parameter constructor produces
\begin{equation}\label{eq:rule}
  D=\left\lceil\frac{\secpar+1+2h}{b}\right\rceil,\qquad
  k=\left\lceil\frac{\secpar}{b}\right\rceil,\qquad
  S=\left\lceil\frac{\secpar}{8}\right\rceil,
\end{equation}
where $b=\lfloor\log_2|\field|\rfloor$ is the conservative bits-per-element.
Here $D$ is the digest width in field elements, and $k$, $S$ are the field-
and byte-salt cardinalities.  Each output is a checkable obligation.  The
constructor either meets it and emits the derived parameters, or it refuses.
A designer can therefore neither quietly under-specify nor silently
over-claim.
\end{rulebox}


The justification is short.  A $\querybudget\le2^{h}$-query collision search succeeds
against a $\digestbits$-bit digest with probability about
$\querybudget^2/2^{\digestbits+1}$, so meeting $2^{-\secpar}$ needs
$\digestbits\ge\secpar+1+2h$ bits, and $D$ is that floor in field elements
(\cref{lem:birthday}).  The salt is attacked one guess at a time, so $k$ and $S$ give it
$\secpar$ bits of entropy in the hash's native unit.

The query budget is where the adversary model enters, and it is the one input the
constructor cannot derive: a prover that grinds a Fiat--Shamir transcript spends its whole
query budget searching for a favourable challenge, which is why the designer declares $h$
rather than the constructor inferring it.

The report follows the rule by property from here.  \Cref{sec:binding} and
\cref{sec:hiding} each carry their own definition, failure, and derived parameter, and
\cref{sec:cost} measures what the derived parameters cost to run.
